% Density revision 2026-10-04: answers below questions; compact related content.
% ৪.১ ওয়েব ডিজাইনের ধারণা
\section{ওয়েব ডিজাইনের ধারণা}

% Composition pass 2026-10-04
কলেজের তথ্য অনলাইনে প্রকাশ করলেই তা সহজে খুঁজে পাওয়া যায় না।

\chFourMotivation{তথ্য দ্রুত খুঁজে পাওয়ার প্রয়োজন $\rightarrow$ ওয়েব ডিজাইন $\rightarrow$ শিরোনাম, মেনু ও বিষয়বস্তু পরিকল্পিতভাবে সাজানো $\rightarrow$ শিক্ষার্থী সহজে নোটিশ ও ফল পায়।}


\begin{learningbox}
এই অংশ শেষে শিক্ষার্থী পারবে:
\begin{itemize}
\item ওয়েব ডিজাইন কী তা সংজ্ঞাসহ ব্যাখ্যা করতে।
\item ওয়েবপেজ, ওয়েবসাইট, ওয়েব সার্ভার ও ওয়েব ব্রাউজারের ধারণা লিখতে।
\item ওয়েব ডিজাইনের উদ্দেশ্য, উপাদান, সুবিধা ও ব্যবহারক্ষেত্র ব্যাখ্যা করতে।
\item HTML (HyperText Markup Language), CSS (Cascading Style Sheets) ও JavaScript-এর ভূমিকা আলাদা করতে।
\item ভালো ওয়েব ডিজাইনের বৈশিষ্ট্য ও পরীক্ষায় গুরুত্বপূর্ণ বিষয় মনে রাখতে।
\end{itemize}
\end{learningbox}

\subsection{ভূমিকা}
ওয়েবসাইটে তথ্য প্রকাশের আগে পেজগুলোর নকশা পরিকল্পনা করতে হয়। কোন তথ্য কোথায় থাকবে, ছবি ও লিংক কীভাবে সাজানো হবে এবং ব্যবহারকারী কোন পথে প্রয়োজনীয় পেজে যাবেন—পরিকল্পনায় এসব ঠিক করা হয়। ব্যক্তি, শিক্ষাপ্রতিষ্ঠান, ব্যবসাপ্রতিষ্ঠান, সংবাদমাধ্যম ও সরকারি সংস্থা এই নকশার সাহায্যে নিজেদের তথ্য ও সেবা উপস্থাপন করে।

এই অধ্যায়ে ওয়েবসাইটের কাঠামো থেকে HTML-এর ব্যবহার পর্যন্ত ধাপে ধাপে আলোচনা করা হয়েছে। লেখার বিন্যাস, তালিকা, হাইপারলিংক, ছবি, সারণি ও ফর্মের উদাহরণ দেখে শিক্ষার্থী ওয়েবপেজ তৈরি, পরীক্ষা ও প্রকাশের কাজ শিখবে।

\subsection{ওয়েব ডিজাইন কী?}

\begin{definitionbox}{ওয়েব ডিজাইন (Web Design)}
ওয়েবসাইট বা ওয়েবপেজের তথ্য, ছবি, লিংক, টেবিল, মেনু ও অন্যান্য উপাদানকে নির্দিষ্ট কাঠামো অনুযায়ী সুন্দর ও ব্যবহারবান্ধবভাবে সাজানোর প্রক্রিয়াকে ওয়েব ডিজাইন বলে।
\end{definitionbox}

ওয়েব ডিজাইনে তথ্য যৌক্তিক ক্রমে সাজানো হয়, যাতে লেখা পড়া এবং প্রয়োজনীয় পেজ খুঁজে পাওয়া সহজ হয়। শুধু বাহ্যিক সৌন্দর্য নয়, মেনু ব্যবহার, তথ্য ডাউনলোড ও সেবা গ্রহণের সুবিধাও নকশার অংশ।

\subsection{ওয়েব ডিজাইনের মূল উদ্দেশ্য}

ওয়েব ডিজাইনের উদ্দেশ্য শুধু সুন্দর পেজ তৈরি করা নয়; বরং তথ্যকে এমনভাবে সাজানো যাতে ব্যবহারকারী দ্রুত, সহজে এবং স্পষ্টভাবে প্রয়োজনীয় তথ্য পায়।

\begin{keypointbox}{প্রধান উদ্দেশ্য}
\par\noindent
\begin{minipage}[t]{.48\linewidth}\vspace{0pt}
\begin{itemize}
\item তথ্য সুন্দর ও সংগঠিতভাবে উপস্থাপন করা।
\item ব্যবহারকারীর জন্য সহজ নেভিগেশন তৈরি করা।
\item ওয়েবসাইটকে আকর্ষণীয় ও ব্যবহারবান্ধব করা।
\item প্রয়োজনীয় তথ্য দ্রুত খুঁজে পাওয়ার ব্যবস্থা করা।
\end{itemize}
\end{minipage}\hfill
\begin{minipage}[t]{.48\linewidth}\vspace{0pt}
\begin{itemize}
\item লেখা, ছবি, লিংক, তালিকা, সারণি ও ফর্মের সঠিক ব্যবহার নিশ্চিত করা।
\item বিভিন্ন ব্রাউজার ও যন্ত্রে ওয়েবপেজ ঠিকভাবে দেখা যায় এমনভাবে তৈরি করা।
\item প্রতিষ্ঠান, পণ্য, সেবা বা ব্যক্তিগত পরিচিতি অনলাইনে প্রকাশ করা।
\item দর্শনার্থী বা ব্যবহারকারী যেন সহজে যোগাযোগ, ডাউনলোড বা সেবা নিতে পারে সে ব্যবস্থা রাখা।
\end{itemize}
\end{minipage}\par

\end{keypointbox}

ওয়েব ডিজাইন বুঝতে ওয়েবপেজ, ওয়েবসাইট ও হোমপেজের সম্পর্ক জানা দরকার। পরের অংশে এগুলোর পরিচয় এবং ক্লায়েন্ট ও সার্ভারের কাজ আলোচনা করা হয়েছে। ওয়েবের ইতিহাস ও ব্রাউজারের ভূমিকা রয়েছে WWW-সংক্রান্ত অংশে।

% Original Bengali explanations after comparison; sources in CHAPTER04_BOOK_COMPARISON.md.
\subsubsection{হাইপারটেক্সট ও হাইপারমিডিয়া}
\textbf{হাইপারটেক্সট (Hypertext)} হলো লিংকের মাধ্যমে সম্পর্কিত নথি বা নথির অংশে যাওয়ার উপযোগী পাঠ। বইয়ের সূচি থেকে অধ্যায়ে যাওয়ার মতো, তবে ওয়েবে লিংক নির্বাচন করেই গন্তব্যে পৌঁছানো যায়। \textbf{হাইপারমিডিয়া (Hypermedia)} এই সংযোগের ধারণাকে লেখা ছাড়াও ছবি, শব্দ ও ভিডিওর মতো মাধ্যমে প্রসারিত করে। যেমন, বিজ্ঞান ক্লাবের পেজের একটি যন্ত্রের ছবি থেকে তার কাজের ব্যাখ্যা বা প্রদর্শনের ভিডিওতে যাওয়া যায়। শুধু একই পেজে নানা মাধ্যম থাকলে সেটি multimedia; লিংক দিয়ে সম্পর্কিত বিষয় অনুসরণের দিকটি hypermedia-তে গুরুত্বপূর্ণ।

\subsubsection{ওয়েব পোর্টাল: সেবার প্রবেশদ্বার}
বিভিন্ন তথ্য, সেবা ও সংশ্লিষ্ট সাইটের লিংক এক জায়গায় সাজিয়ে দেওয়া বিশেষ ধরনের ওয়েবসাইটকে \textbf{ওয়েব পোর্টাল (Web Portal)} বলা হয়। একটি শিক্ষার্থী পোর্টালে নোটিশ, কোর্স, পরীক্ষার রুটিন ও ফলাফলের প্রবেশপথ থাকতে পারে। প্রতিটি সেবার ফাইল একই সার্ভারে থাকতেই হবে না। \textbf{পার্থক্য:} homepage একটি সাইটের প্রধান পেজ; portal একটি সাইটের উদ্দেশ্য বা সেবার ধরন। তাই প্রতিটি homepage-কে portal বলা ঠিক নয়।

কলেজের রুটিন, ফি ও ফলাফলের সেবায় এক জায়গা থেকে যাওয়ার ব্যবস্থা থাকলে সেই সাইট পোর্টালের উদাহরণ। শুধু সাধারণ পরিচিতির পেজে এই বৈশিষ্ট্য নাও থাকতে পারে।


\subsection{নকশা থেকে পর্দায় ওয়েবপেজ}
নকশা অনুযায়ী HTML-এ পেজের কাঠামো, CSS-এ তার দৃশ্যমান বিন্যাস এবং প্রয়োজনে JavaScript-এ ব্যবহারকারীর কাজের প্রতিক্রিয়া তৈরি করা হয়। প্রকাশের পর ব্রাউজার সার্ভার থেকে এসব ফাইল নিয়ে পেজ দেখায়। অনুরোধ ও সাড়া আদান-প্রদানের বিস্তারিত ধাপ ক্লায়েন্ট–সার্ভার ও ওয়েব ঠিকানা অংশে রয়েছে।

\subsection{ওয়েব ডিজাইনে HTML-এর ভূমিকা}
HTML দিয়ে শিরোনাম, অনুচ্ছেদ, ছবি, তালিকা, সারণি, ফর্ম ও লিংকের অবস্থান এবং পারস্পরিক গঠন লেখা হয়। CSS ও JavaScript কীভাবে এই কাঠামোর সঙ্গে কাজ করে, তা HTML-এর মৌলিক বিষয় অংশে ব্যাখ্যা করা হয়েছে।

\subsection{ওয়েব ডিজাইনের জন্য প্রয়োজনীয় উপাদান}

একটি ভালো ওয়েবপেজ ডিজাইনের জন্য সাধারণত নিচের বিষয়গুলো দরকার হয়:

\begin{center}
\begin{tabular}{|p{0.28\textwidth}|p{0.54\textwidth}|}
\hline
\textbf{উপাদান} & \textbf{ব্যাখ্যা} \\
\hline
বিষয়বস্তু & লেখা, তথ্য, ছবি, ভিডিও, সারণি, নোটিশ, ফল ইত্যাদি \\
\hline
কাঠামো & শিরোনাম, অনুচ্ছেদ, অংশ, তালিকা, সারণি, ফর্মের বিন্যাস \\
\hline
নেভিগেশন & মেনু, হাইপারলিংক, বোতাম, back/home লিংক \\
\hline
দৃশ্যমান নকশা & রং, ফন্ট, ফাঁক, ছবি আকার, বিন্যাস \\
\hline
ব্যবহারের সহজতা & ব্যবহারকারী যেন সহজে বুঝতে ও ব্যবহার করতে পারে \\
\hline
সবার ব্যবহারযোগ্যতা & বিভিন্ন ব্যবহারকারী যেন পেজ ব্যবহার করতে পারে \\
\hline
কাজের গতি & পেজ দ্রুত লোড হওয়া \\
\hline
সামঞ্জস্য & বিভিন্ন ব্রাউজার ও যন্ত্রে ঠিকভাবে দেখা যাওয়া \\
\hline
\end{tabular}
\end{center}

% Related compact discussions; source reading order is retained.
\par\noindent
\begin{minipage}[t]{.48\linewidth}\vspace{0pt}
\subsection{ওয়েব ডিজাইনের সুবিধা}

\begin{enumerate}
\item \textbf{তথ্য প্রকাশ:} পরিকল্পিত ওয়েবসাইটে ব্যক্তি ও প্রতিষ্ঠানের তথ্য সহজে ও দ্রুত প্রকাশ করা যায়।
\item \textbf{তথ্যে প্রবেশ:} সার্ভার ও সংযোগ সচল থাকলে শিক্ষার্থী, অভিভাবক ও গ্রাহক বিভিন্ন স্থান থেকে ২৪ ঘণ্টা প্রয়োজনীয় তথ্য দেখতে পারেন।
\item প্রতিষ্ঠান বা ব্যবসার পরিচিতি বাড়ে।
\item অনলাইন সেবা দেওয়া যায়।
\item শিক্ষার্থীরা অনলাইন ক্লাস, নোট, ফল, রুটিন ও নোটিশ পেতে পারে।
\item গ্রাহক বা দর্শনার্থী সহজে যোগাযোগ করতে পারে।
\item কম খরচে বিশ্বব্যাপী প্রচার করা যায়।
\item প্রয়োজনীয় ফাইল, PDF, ফর্ম বা নোটিশ ডাউনলোড করা যায়।
\end{enumerate}
\end{minipage}\hfill
\begin{minipage}[t]{.48\linewidth}\vspace{0pt}
\subsection{ওয়েব ডিজাইনের ব্যবহারক্ষেত্র}

\begin{enumerate}
\item \textbf{শিক্ষা প্রতিষ্ঠান:} নোটিশ, ফল, ভর্তি, রুটিন, শিক্ষকদের তথ্য।
\item \textbf{ব্যবসা প্রতিষ্ঠান:} পণ্য, সেবা, অর্ডার, অর্থ পরিশোধের তথ্য।
\item \textbf{সংবাদমাধ্যম:} অনলাইন সংবাদ পোর্টাল।
\item \textbf{সরকারি সেবা:} নাগরিক সেবা, আবেদন, তথ্য প্রকাশ।
\item \textbf{স্বাস্থ্যসেবা:} সাক্ষাতের সময় নেওয়া, চিকিৎসকের তথ্য, স্বাস্থ্যসংক্রান্ত নোটিশ।
\item \textbf{ব্যক্তিগত পোর্টফোলিও:} নিজের কাজ, দক্ষতা, প্রকল্প দেখানো।
\item \textbf{ই-কমার্স:} অনলাইন কেনাকাটা, পণ্যের তালিকা, কেনাকাটার ঝুড়ি।
\item \textbf{ব্যাংকিং:} অনলাইন ব্যাংকিং তথ্য ও সেবা নোটিশ।
\end{enumerate}
\end{minipage}\par\smallskip
\subsection{ভালো ওয়েব ডিজাইনের বৈশিষ্ট্য}

একটি ভালো ওয়েবসাইট শুধু দেখতে সুন্দর হলেই যথেষ্ট নয়; সেটি ব্যবহার করাও সহজ হতে হবে।

\begin{keypointbox}{ভালো ওয়েব ডিজাইনের বৈশিষ্ট্য}
\par\noindent
\begin{minipage}[t]{.48\linewidth}\vspace{0pt}
\begin{enumerate}
\item পরিষ্কার ও সহজ বিন্যাস।
\item সহজ নেভিগেশন মেনু।
\item পাঠযোগ্য ফন্ট ও রং।
\item দ্রুত পেজ লোড হওয়া।
\item প্রয়োজনীয় ছবি ও বিষয়বস্তু।
\item মোবাইলে ব্যবহারযোগ্য নকশা।
\end{enumerate}
\end{minipage}\hfill
\begin{minipage}[t]{.48\linewidth}\vspace{0pt}
\begin{enumerate}[start=7]
\item অকার্যকর লিংক না থাকা।
\item ব্যবহারবান্ধব কাঠামো।
\item তথ্য নিয়মিত হালনাগাদ করা।
\item ব্রাউজারের সঙ্গে সামঞ্জস্য থাকা।
\item অপ্রয়োজনীয় অ্যানিমেশন বা অতিরিক্ত নকশা এড়িয়ে চলা।
\end{enumerate}
\end{minipage}\par

\end{keypointbox}

% Original applied explanation; see Chapter04/SOURCES.md.
\subsubsection{নোটিশ খুঁজতে গিয়ে ডিজাইনের প্রয়োজন বোঝা}
একজন অভিভাবক কলেজের ভর্তি-তারিখ জানতে চান। হোমপেজে বড় ছবি থাকলেও ভর্তি-নোটিশের লিংক অস্পষ্ট হলে তাঁর কাজ হয় না। মেনুতে \textbf{ভর্তি তথ্য}, নোটিশে স্পষ্ট তারিখ এবং প্রয়োজনীয় ফাইলের অর্থপূর্ণ লিংক দিলে তিনি সহজে তথ্য পান।

ডিজাইনের মান তাই দর্শকের কাজ দিয়ে বিচার করবে: প্রয়োজনীয় তথ্য পাওয়া গেল কি, লেখা পড়া গেল কি এবং পরের ধাপ বোঝা গেল কি? এই তিনটি প্রশ্ন দিয়ে নিজের কলেজ, দোকান বা গ্রন্থাগারের পেজ পরীক্ষা করা যায়।


\subsection{ওয়েব ডিজাইন বনাম ওয়েব ডেভেলপমেন্ট}

\begin{center}
\begin{tabular}{|p{0.20\textwidth}|p{0.32\textwidth}|p{0.32\textwidth}|}
\hline
\textbf{বিষয়} & \textbf{ওয়েব ডিজাইন} & \textbf{ওয়েব ডেভেলপমেন্ট} \\
\hline
কাজ & ওয়েবপেজের দৃশ্যমান রূপ ও বিন্যাস তৈরি & ওয়েবসাইটের কার্যক্ষমতা তৈরি \\
\hline
গুরুত্ব & রং, ফন্ট, বিন্যাস, ছবি, ব্যবহারকারীর অভিজ্ঞতা & কোড, ডেটাবেজ, সার্ভার, যুক্তি \\
\hline
উপকরণ & HTML, CSS, নকশা তৈরির উপকরণ & HTML, CSS, JavaScript, PHP (PHP: Hypertext Preprocessor), ডেটাবেজ \\
\hline
উদাহরণ & হোমপেজ সুন্দরভাবে সাজানো & লগইন ব্যবস্থা বা ফল অনুসন্ধান ব্যবস্থা তৈরি \\
\hline
\end{tabular}
\end{center}



\subsection{ওয়েবসাইট তৈরির ধাপ}

\begin{enumerate}
\item ওয়েবসাইটের উদ্দেশ্য ও লক্ষ্য ব্যবহারকারী নির্ধারণ করা।
\item কোন কোন পেজ থাকবে তা ঠিক করা।
\item বিষয়বস্তু, ছবি, ফাইল ও লিংক সংগ্রহ করা।
\item সাইটম্যাপ বা পেজের কাঠামো তৈরি করা।
\item খসড়া বিন্যাস বা ওয়্যারফ্রেম তৈরি করা।
\item HTML দিয়ে ওয়েবপেজের কাঠামো তৈরি করা।
\item প্রয়োজনে CSS দিয়ে পেজ সুন্দর করা।
\item ব্রাউজারে ফল পরীক্ষা করা।
\item লিংক, ছবি, সারণি ও ফর্ম ঠিক আছে কিনা পরীক্ষা করা।
\item হোস্টিং সার্ভারে ওয়েবসাইট প্রকাশ করা।
\end{enumerate}

ওয়েবসাইটের কাজ অনুযায়ী প্রযুক্তি বেছে নিতে হয়। পরিচিতি প্রকাশ ও রোলভিত্তিক ফল অনুসন্ধানের চাহিদা এক নয়। স্ট্যাটিক, ডাইনামিক, লোকাল ও রিমোট সাইটের পার্থক্য ওয়েবসাইটের প্রকারভেদ অংশে আলোচনা করা হয়েছে।

\subsection{একটি সাধারণ HTML কাঠামো}

নিচের কোডটি একটি সাধারণ ওয়েবপেজের মৌলিক কাঠামো দেখায়।

% Split revision: code and its explanation share one row.
\par\noindent
\begin{minipage}[t]{.54\linewidth}\vspace{0pt}\small
\textbf{\color{chFourTeal}কোড}\par
\begin{verbatim}
<!DOCTYPE html>
<html>
<head>
  <meta charset="UTF-8">
  <title>Example College</title>
</head>
<body>
  <h1>Example College</h1>
  <p>Welcome to our college website.</p>
  <a href="notice.html">Notice</a>
  <br>
  <img src="images/college.jpg" alt="College" width="300">
</body>
</html>
\end{verbatim}
\par\smallskip
\end{minipage}\hfill
\begin{minipage}[t]{.42\linewidth}\vspace{0pt}\small
\textbf{\color{chFourTeal}ফল সম্পর্কে ধারণা}\par
ব্রাউজারের ট্যাবে \texttt{Example College} লেখা দেখা যাবে। পেজে থাকবে কলেজের শিরোনাম, স্বাগতবার্তা, Notice লিংক ও কলেজের ছবি। \texttt{images} ফোল্ডারে \texttt{college.jpg} ফাইলটি থাকলে ছবিটি দেখা যাবে।
\par\smallskip
\end{minipage}\par


\begin{keypointbox}{কোডের ব্যাখ্যা}
\par\noindent
\begin{minipage}[t]{.48\linewidth}\vspace{0pt}
\begin{itemize}
\item \texttt{<!DOCTYPE html>} আধুনিক HTML-এর ডকুমেন্ট ঘোষণা; এটি ব্রাউজারকে মানসম্মত মোডে পেজ দেখাতে সাহায্য করে।
\item \texttt{<title>} ব্রাউজারের ট্যাবের শিরোনাম নির্ধারণ করে।
\item \texttt{<h1>} সর্বোচ্চ স্তরের শিরোনাম তৈরি করে; সাধারণ বিন্যাসে এটি বড় আকারে দেখা যায়।
\end{itemize}
\end{minipage}\hfill
\begin{minipage}[t]{.48\linewidth}\vspace{0pt}
\begin{itemize}
\item \texttt{<p>} অনুচ্ছেদ তৈরি করে।
\item \texttt{<a href="notice.html">} অভ্যন্তরীণ হাইপারলিংক তৈরি করে।
\item \texttt{<img>} ওয়েবপেজে ছবি দেখায়।
\end{itemize}
\end{minipage}\par

\end{keypointbox}

\Needspace{10\baselineskip}
\subsection{সংক্ষিপ্ত অনুশীলন}

\chFourReviewQuestion{প্রশ্ন: ওয়েব ডিজাইন কী?}

\chFourReviewQuestion{প্রশ্ন: ওয়েব ডিজাইন কেন গুরুত্বপূর্ণ?}


\subsection{বহুনির্বাচনি অনুশীলন}

\begin{multicols}{2}
\begin{enumerate}[leftmargin=*,itemsep=5pt]
\item \begin{minipage}[t]{\linewidth}\vspace{0pt}
ওয়েবপেজ তৈরিতে কোন ভাষা বেশি ব্যবহৃত হয়?
\par ক) Python
\par খ) HTML
\par গ) COBOL
\par ঘ) FORTRAN\\
\chFourMCQAnswer{খ) HTML}
\end{minipage}

\item \begin{minipage}[t]{\linewidth}\vspace{0pt}
একই ডোমেইনের অধীনে একাধিক ওয়েবপেজের সমষ্টিকে কী বলে?
\par ক) Web browser
\par খ) Website
\par গ) Search engine
\par ঘ) Web server\\
\chFourMCQAnswer{খ) Website}
\end{minipage}

\item \begin{minipage}[t]{\linewidth}\vspace{0pt}
ওয়েবসাইট দেখার জন্য কোন সফটওয়্যার ব্যবহার করা হয়?
\par ক) Compiler
\par খ) Browser
\par গ) Assembler
\par ঘ) Spreadsheet\\
\chFourMCQAnswer{খ) Browser}
\end{minipage}

\item \begin{minipage}[t]{\linewidth}\vspace{0pt}
HTML-এর পূর্ণরূপ কী?
\par ক) HighText Machine Language
\par খ) HyperText Markup Language
\par গ) Hyper Tool Mark Language
\par ঘ) Home Text Markup Language\\
\chFourMCQAnswer{খ) HyperText Markup Language}
\end{minipage}
\end{enumerate}
\end{multicols}

একটি কলেজের সাইটে Home, About College, Teachers, Notice, Admission, Result ও Contact পেজ রাখা যেতে পারে। প্রতিটি পেজে তার বিষয় অনুযায়ী তথ্য থাকে।

\subsection{কাজ}

\begin{practicebox}
তোমার কলেজের জন্য একটি ওয়েবসাইট কল্পনা করো। সেখানে কোন কোন পেজ থাকবে লিখো।

\end{practicebox}

